\documentclass[letterpaper, 12pt, twoside]{article}
\usepackage{amsmath}
\usepackage{graphicx}
\usepackage{xcolor}
\usepackage{enumitem}
\usepackage[utf8]{inputenc}
\usepackage{{graphicx}}

%Aquí inicia la portada del documento
\title{\Huge\item\color{red}\textbf {VITÁCORA DEL DÍA NÚMERO 01 DEL CURSO}}
\author{Mariana Yasmin Martínez García}
\date{13/01/2019}

\begin{document}
	\maketitle
	\begin{figure}
		\centering
		%\includegraphics{/home/mariana/Documentos/TALLER-DE-HERRAMIENTAS-COMPUTACIONALES/Clases/Latex/Imágenes/1}
		\caption{}
		\label{fig:1}
	\end{figure}
	
	
	\newpage
	
	\title{\huge\textbf{Vitácora del día 7 de enero\\}} \\
	El primer día del curso se nos dió de forma sencilla la forma de evaluación al igual que la introducción de lo que veríamos en el curso.
	Se vió que existe un sistema operativo llamado \textbf{Linux} que es gratis y estable, y hay varias distribuciones distintas de él, algunas de ellas y las más conocidas son:
	\begin{enumerate}
		\item\textbf{Fedora}
		\item\textbf{Ubuntu}
	\end{enumerate} 
    Se nos dijo que hay varias variantes de estas distribuciones y que podemos colocar las de manera gratuita en nuestro equipo. Y se nos pidió que intentáramos descargarlo ya que lo íbamos a estar ocupando por el resto el curso.\\ \\
   
También se nos enseñó algunos comandos, como son: 
	\begin{enumerate}
		\item \textbf{\large pwd }: sirve para visualizar el directorio en el que nos encontramos.
		\item \textbf{\large ls }: este te muestra los archivos que hay en algún directorio.
		\item \textbf{\large ls -a }: este te muestra TODOS los archivos que hay en algún directorio, incluyendo los archivos ocultos.
		\item \textbf{\large ls -l }: te en lista las propiedades y permisos de algún archivo.
		\item \textbf{\large chmod }: con este puedes modificar los permisos de algún archivo.
		\item \textbf{\large mkdir \textit{directorio} }: crea un nuevo directorio.
		\item \textbf{\large cd \textit{directorio}}: se escribe el comando seguido del nombre de un directorio y te dirigirá a este.
		\item \textbf{\large cd .. :} te dirige al directorio donde estabas antes.
		\item \textbf{\large man \textit{comando} }: te muestra el manual del comando que escribiste.	
	\end{enumerate}
    Y ya que esta primer clase fue en un salón que no contaba con computadoras, se dejó de tarea probar que todos estos comandos sí funcionaran y practicar con ellos.

	
	
	
\end{document}